% !TeX root = ../main.tex
\begin{abstract}
In multi-event Known-Item Search (KIS) at the Video Browser Showdown,
complex narrative queries typically describe sequences of temporally ordered
visual events. While large language models (LLMs) can decompose such narratives
into structured sub-queries, errors in event boundaries or temporal sequencing
frequently cascade into downstream retrieval failures. Such errors remain
opaque when query parsing and temporal alignment operate as internal,
black-box processes. We present SHI (Structured Hypothesis Inspection), an
interactive video retrieval system that externalizes intermediate
representations into explicit interaction objects across two complementary
stages. In the pre-retrieval stage, the Query Hypothesis exposes the
LLM decomposition as an editable ordered sequence, enabling users to verify
semantics, reorder events, and bind reference images before querying. In the
post-retrieval stage, EventTrail visualizes the temporal alignment of candidate
videos as coherent ordered paths rather than isolated keyframes. Users can
anchor valid occurrences or exclude false-positive intervals via keep, use,
and reject actions, triggering dynamic-programming path re-decoding over
retained multimodal evidence snapshots without repeating corpus-level
retrieval. In an exploratory evaluation on multi-event KIS tasks, SHI achieved
a 90\% valid timestamp hit rate compared to 60\% for a standard baseline,
demonstrating that structured hypothesis inspection effectively recovers from
early-stage cascade errors.

\keywords{
Interactive video retrieval
\and Video Browser Showdown
\and Temporal alignment
\and Human-in-the-loop
\and Multi-event search
}
\end{abstract}